\documentclass[10pt,a4paper]{amsart}
\usepackage{amsmath,amssymb,amsthm,graphicx}
% Retained minimal preamble for the Stacks Project chapter more-algebra.tex (More on Algebra).
% The chapter file pinned-original-source.tex carries only \input{preamble} on its first line, so no
% preamble text lives inside the chapter itself.  The Stacks Project uses one single shared preamble for
% every chapter; the snapshot stages it identically as sources/stacks-chow/preamble.tex and
% sources/stacks-etale-cohomology/preamble.tex.  The block below is copied mechanically, line for line,
% from sources/stacks-chow/preamble.tex lines 14-18 (the Xy-pic setup) and 154-217 (the native theorem
% declarations together with the mathematical macros).  Nothing was added, removed or redefined; the
% wrapper class and the amsmath/amssymb/amsthm/graphicx base packages are supplied by the assembler.
% Chapter: more-algebra.tex -- More on Algebra.
% For commutative diagrams we use Xy-pic
\usepackage[all]{xy}

% We use 2cell for 2-commutative diagrams.
\xyoption{2cell}

% Theorem environments.
%
\theoremstyle{plain}
\newtheorem{theorem}[subsection]{Theorem}
\newtheorem{proposition}[subsection]{Proposition}
\newtheorem{lemma}[subsection]{Lemma}

\theoremstyle{definition}
\newtheorem{definition}[subsection]{Definition}
\newtheorem{example}[subsection]{Example}
\newtheorem{exercise}[subsection]{Exercise}
\newtheorem{situation}[subsection]{Situation}

\theoremstyle{remark}
\newtheorem{remark}[subsection]{Remark}
\newtheorem{remarks}[subsection]{Remarks}

\numberwithin{equation}{subsection}

% Macros
%
\def\lim{\mathop{\mathrm{lim}}\nolimits}
\def\colim{\mathop{\mathrm{colim}}\nolimits}
\def\Spec{\mathop{\mathrm{Spec}}}
\def\Hom{\mathop{\mathrm{Hom}}\nolimits}
\def\Ext{\mathop{\mathrm{Ext}}\nolimits}
\def\SheafHom{\mathop{\mathcal{H}\!\mathit{om}}\nolimits}
\def\SheafExt{\mathop{\mathcal{E}\!\mathit{xt}}\nolimits}
\def\Sch{\mathit{Sch}}
\def\Mor{\mathop{\mathrm{Mor}}\nolimits}
\def\Ob{\mathop{\mathrm{Ob}}\nolimits}
\def\Sh{\mathop{\mathit{Sh}}\nolimits}
\def\NL{\mathop{N\!L}\nolimits}
\def\CH{\mathop{\mathrm{CH}}\nolimits}
\def\proetale{{pro\text{-}\acute{e}tale}}
\def\etale{{\acute{e}tale}}
\def\QCoh{\mathit{QCoh}}
\def\Ker{\mathop{\mathrm{Ker}}}
\def\Im{\mathop{\mathrm{Im}}}
\def\Coker{\mathop{\mathrm{Coker}}}
\def\Coim{\mathop{\mathrm{Coim}}}

% Boxtimes
%
\DeclareMathSymbol{\boxtimes}{\mathbin}{AMSa}{"02}

%
% Macros for moduli stacks/spaces
%
\def\QCohstack{\mathcal{QC}\!\mathit{oh}}
\def\Cohstack{\mathcal{C}\!\mathit{oh}}
\def\Spacesstack{\mathcal{S}\!\mathit{paces}}
\def\Quotfunctor{\mathrm{Quot}}
\def\Hilbfunctor{\mathrm{Hilb}}
\def\Curvesstack{\mathcal{C}\!\mathit{urves}}
\def\Polarizedstack{\mathcal{P}\!\mathit{olarized}}
\def\Complexesstack{\mathcal{C}\!\mathit{omplexes}}
% \Pic is the operator that assigns to X its picard group, usage \Pic(X)
% \Picardstack_{X/B} denotes the Picard stack of X over B
% \Picardfunctor_{X/B} denotes the Picard functor of X over B
\def\Pic{\mathop{\mathrm{Pic}}\nolimits}
\def\Picardstack{\mathcal{P}\!\mathit{ic}}
\def\Picardfunctor{\mathrm{Pic}}
\def\Deformationcategory{\mathcal{D}\!\mathit{ef}}

% --- appended 2026-10-12 (Chief): the project own definitions of the reference/slogan/history
% display environments, copied mechanically from the shared Stacks preamble lines 8-12
% (\usepackage{verbatim} plus the three \newenvironment lines, which the chapter preamble keeps
% verbatim from the project preamble). Nothing added, removed or redefined: these environments
% carry no mathematical content and the project itself typesets them as comments. Without these
% three lines a counted statement whose own text uses them fails with Environment slogan undefined.
\usepackage{verbatim}
\newenvironment{reference}{\comment}{\endcomment}
\newenvironment{slogan}{\comment}{\endcomment}
\newenvironment{history}{\comment}{\endcomment}
\usepackage{cleveref}
\crefname{originallemma}{original lemma}{original lemmas}
\Crefname{originallemma}{Original lemma}{Original lemmas}
\makeatletter
\newlabel{derived-lemma-morphisms-from-projective-complex}{{064B}{source}{}{sourcepointer-1}{}}
\newlabel{derived-lemma-morphisms-from-projective-complex@cref}{{[originallemma][0][]064B}{[1][0][]source}{}{}{}}
\makeatother
\crefname{originallemma}{original lemma}{original lemmas}
\Crefname{originallemma}{Original lemma}{Original lemmas}
\makeatletter
\newlabel{lemma-canonical-element-well-defined}{{36985}{source}{}{sourcepointer-2}{}}
\newlabel{lemma-canonical-element-well-defined@cref}{{[originallemma][0][]36985}{[1][0][]source}{}{}{}}
\makeatother
\crefname{originallemma}{original lemma}{original lemmas}
\Crefname{originallemma}{Original lemma}{Original lemmas}
\makeatletter
\newlabel{lemma-compose-surjections}{{37108}{source}{}{sourcepointer-3}{}}
\newlabel{lemma-compose-surjections@cref}{{[originallemma][0][]37108}{[1][0][]source}{}{}{}}
\makeatother
\crefname{originallemma}{original lemma}{original lemmas}
\Crefname{originallemma}{Original lemma}{Original lemmas}
\makeatletter
\newlabel{lemma-homotopic-surjections}{{37016}{source}{}{sourcepointer-4}{}}
\newlabel{lemma-homotopic-surjections@cref}{{[originallemma][0][]37016}{[1][0][]source}{}{}{}}
\makeatother
\crefname{originallemma}{original lemma}{original lemmas}
\Crefname{originallemma}{Original lemma}{Original lemmas}
\makeatletter
\newlabel{lemma-perfect}{{18889}{source}{}{sourcepointer-5}{}}
\newlabel{lemma-perfect@cref}{{[originallemma][0][]18889}{[1][0][]source}{}{}{}}
\makeatother
\title{The determinant of a perfect complex of tor amplitude in two degrees is a functor to invertible modules, with a canonical element carrying isomorphisms}
\author{The Stacks Project Authors}
\date{stacks-project:more-algebra.tex:local-snapshot-20261012; GNU FDL 1.2 or later}
\begin{document}
\maketitle
\setcounter{subsection}{1}
\noindent\textit{Editorial note.} All delivered mathematical text is the author's own, copied from the pinned Stacks Project chapter file more-algebra.tex; no statement or proof step was rewritten, reordered or removed.\par\medskip
\noindent\textit{Editorial note.} The chapter file pinned-original-source.tex carries only \textbackslash{}input\{preamble\} on its first line, so no preamble text lives inside the chapter. The retained file sources/stacks-more-algebra/retained-preamble.tex is a mechanical, line-for-line copy of the project's single shared preamble as staged in the snapshot; nothing was added, removed or redefined and no package was added.\par\medskip
\noindent\textit{Editorial note.} The source URL is the source-level Stacks Project URL, as required when more than one card is produced from one source; the per-claim tag is carried by claim\_key.\par\medskip
\noindent\textit{Editorial note.} The counted statement span contains the author's two large displayed categories, delivered verbatim. The five references inside the counted proof span are recorded as non-load-bearing navigation pointers.\par\medskip
\noindent\textit{Allowed standard prerequisites (editorial).}\begin{itemize}
\item The determinant of a two-term complex of finite projective modules as det(L\textasciicircum{}0) tensor det(L\textasciicircum{}\{-1\})\textasciicircum{}\{tensor -1\}. used at proof lines 37160-37162. Is the object the functor is built from.
\item The three conditions on the maps of a good diagram established earlier in the section. used at proof lines 37173-37181 and 37198-37200. Are what make det(c\textasciicircum{}bullet) o det(b\textasciicircum{}bullet)\textasciicircum{}\{-1\} a sensible definition.
\item The homotopic-surjections lemma comparing the determinants of two maps that agree up to homotopy. invoked at proof line 37210. Gives independence of det(a\textasciicircum{}bullet) from the chosen good diagram in the identity case.
\item Representation of morphisms from a bounded complex of finite projective modules by homotopy classes of maps of complexes. invoked at proof lines 37156-37159. Lets the construction be carried out on good complexes rather than abstractly.
\end{itemize}\medskip
\ifdefined\hypertarget\hypertarget{sourcepointer-1}{}\fi
\noindent\textit{Source reference.} Original lemma 064B: focused-seeds/AGP/sources/stacks-snapshot-20261011/tags.txt, source lines 7860--7860. Cross-chapter navigation pointer into the Stacks Project; the standard fact it supplies is applied inside the counted proof span and the argument resting on it stays there.\par\medskip
\ifdefined\hypertarget\hypertarget{sourcepointer-2}{}\fi
\noindent\textit{Source reference.} Original lemma 36985: sources/stacks-more-algebra/pinned-original-source.tex, source lines 36985--36985. Same-chapter navigation pointer; the statement it names is standard for the section and the argument that consumes it stays inside the counted proof span.\par\medskip
\ifdefined\hypertarget\hypertarget{sourcepointer-3}{}\fi
\noindent\textit{Source reference.} Original lemma 37108: sources/stacks-more-algebra/pinned-original-source.tex, source lines 37108--37108. Same-chapter navigation pointer; the statement it names is standard for the section and the argument that consumes it stays inside the counted proof span.\par\medskip
\ifdefined\hypertarget\hypertarget{sourcepointer-4}{}\fi
\noindent\textit{Source reference.} Original lemma 37016: sources/stacks-more-algebra/pinned-original-source.tex, source lines 37016--37016. Same-chapter navigation pointer; the statement it names is standard for the section and the argument that consumes it stays inside the counted proof span.\par\medskip
\ifdefined\hypertarget\hypertarget{sourcepointer-5}{}\fi
\noindent\textit{Source reference.} Original lemma 18889: sources/stacks-more-algebra/pinned-original-source.tex, source lines 18889--18889. Same-chapter navigation pointer; the statement it names is standard for the section and the argument that consumes it stays inside the counted proof span.\par\medskip
\begin{lemma}
\label{lemma-determinant-two-term-complexes}
Let $R$ be a ring. The constructions above determine a functor
$$
\det :
\left\{
\begin{matrix}
\text{category of perfect complexes} \\
\text{with tor amplitude in }[-1, 0] \\
\text{morphisms are isomorphisms}
\end{matrix}
\right\}
\longrightarrow
\left\{
\begin{matrix}
\text{category of invertible modules} \\
\text{morphisms are isomorphisms}
\end{matrix}
\right\}
$$
Moreover, given a rank $0$ perfect object $L$ of $D(R)$ with
tor-amplitude in $[-1, 0]$ there is a canonical element
$\delta(L) \in \det(L)$ such that for any isomorphism
$a : L \to K$ in $D(R)$ we have $\det(a)(\delta(L)) = \delta(K)$.
\end{lemma}

\begin{proof}
By Lemma \ref{lemma-perfect} every object of the source category
may be represented by a complex
$$
L^\bullet = \ldots \to 0 \to L^{-1} \to L^0 \to 0 \to \ldots
$$
with $L^{-1}$ and $L^0$ finite projective $R$-modules. Let us temporarily
call a complex of this type good. By 
Derived Categories, Lemma \ref{derived-lemma-morphisms-from-projective-complex}
morphisms between good complexes in the derived category are
homotopy classes of maps of complexes. Thus we may work with good complexes
and we can use the determinant
$\det(L^\bullet) = \det(L^0) \otimes \det(L^{-1})^{\otimes -1}$
we investigated above.

\medskip\noindent
Let $a^\bullet : L^\bullet \to K^\bullet$ be a morphism of good complexes
which is an isomorphism in $D(R)$, i.e., a quasi-isomorphism. We say that
$$
\xymatrix{
L^\bullet \ar[rr]_{a^\bullet} & & K^\bullet \\
& M^\bullet \ar[lu]^{b^\bullet} \ar[ru]_{c^\bullet}
}
$$
is a good diagram if it commutes up to homotopy and $b^\bullet$
and $c^\bullet$ satisfy conditions (1), (2), (3) above. Whenever
we have such a diagram it makes sense to define
$$
\det(a^\bullet) = \det(c^\bullet) \circ \det(b^\bullet)^{-1}
$$
where $\det(c^\bullet)$ and $\det(b^\bullet)$ are the isomorphisms
constructed in the text above. We will show that good diagrams always
exist and that the resulting map $\det(a^\bullet)$
is independent of the choice of good diagram.

\medskip\noindent
Existence of good diagrams for a quasi-isomorphism
$a^\bullet : L^\bullet \to K^\bullet$ of good complexes.
Choose a
surjection $p : R^{\oplus n} \to K^{-1}$. Then we can consider the new
good complex
$$
M^\bullet = \ldots \to 0 \to
L^{-1} \oplus R^{\oplus n} \xrightarrow{d \oplus 1}
L^0 \oplus R^{\oplus n} \to 0 \to \ldots
$$
with the projection map $b^\bullet : M^\bullet \to L^\bullet$ and the map
$c^\bullet : M^\bullet \to K^\bullet$ using $a^{-1} \oplus p$ in degree $-1$
and using $a^0 \oplus d \circ p$ in degree $0$. The maps
$b^\bullet : M^\bullet \to L^\bullet$ and
$c^\bullet : M^\bullet \to K^\bullet$
satisfy conditions (1), (2), (3) above and we get a good diagram.

\medskip\noindent
Suppose that we have a good diagram
$$
\xymatrix{
L^\bullet \ar[rr]_{\text{id}^\bullet} & & L^\bullet \\
& M^\bullet \ar[lu]^{b^\bullet} \ar[ru]_{c^\bullet}
}
$$
Then by Lemma \ref{lemma-homotopic-surjections} we see that
$\det(c^\bullet) = \det(b^\bullet)$. Thus we see that
$\det(\text{id}^\bullet) = \text{id}$ is independent of the choice
of good diagram.

\medskip\noindent
Before we prove independence in general, we think about
composition. Suppose we have quasi-isomorphisms
$L_1^\bullet \to L_2^\bullet$ and $L_2^\bullet \to L_3^\bullet$
of good complexes and good diagrams
$$
\vcenter{
\xymatrix{
L_1^\bullet \ar[rr] & &
L_2^\bullet \\
& M_{12}^\bullet \ar[lu] \ar[ru]
}
}
\quad\text{and}\quad
\vcenter{
\xymatrix{
L_2^\bullet \ar[rr] & &
L_3^\bullet \\
& M_{23}^\bullet \ar[lu] \ar[ru]
}
}
$$
We can extend this to a diagram
$$
\xymatrix{
L_1^\bullet \ar[rr] & &
L_2^\bullet \ar[rr] & &
L_3^\bullet \\
& M_{12}^\bullet \ar[lu] \ar[ru]
& & M_{23}^\bullet \ar[lu] \ar[ru] \\
& &
M_{123}^\bullet \ar[lu] \ar[ru]
}
$$
where $M_{123}^\bullet \to M_{12}^\bullet$ and
$M_{123}^\bullet \to M_{23}^\bullet$ have properties (1), (2), (3)
and the square in the diagram commutes: we can just take
$M_{123}^n = M_{12}^n \times_{L_2^n} M_{23}^n$.
Then Lemma \ref{lemma-compose-surjections} shows that
$$
\xymatrix{
\det(L_2^\bullet) &
\det(M_{23}^\bullet) \ar[l] \\
\det(M_{12}^\bullet) \ar[u] &
\det(M_{123}^\bullet) \ar[l] \ar[u]
}
$$
commutes. A diagram chase shows that the composition
$\det(L_1^\bullet) \to \det(L_2^\bullet) \to \det(L_3^\bullet)$
of the maps associated to the two good diagrams using
$M_{12}^\bullet$ and $M_{23}^\bullet$
is equal to the map associated to the good diagram
$$
\xymatrix{
L_1^\bullet \ar[rr] & & L_3^\bullet \\
& M_{123}^\bullet \ar[lu] \ar[ru]
}
$$
Thus if we can show that these maps are independent of choices,
then the composition law is satisfied too and we obtain our functor.

\medskip\noindent
Independence. Let a quasi-isomorphism
$a^\bullet : L^\bullet \to K^\bullet$ of good complexes be given.
Choose an inverse quasi-isomorphism $b^\bullet : K^\bullet \to L^\bullet$.
Setting $L_1^\bullet = L$, $L_2^\bullet = K^\bullet$ and
$L_3^\bullet = L^\bullet$ may fix our choice of good diagram for
$b^\bullet$ and consider varying good diagrams for $a^\bullet$.
Then the result of the previous paragraphs is that no matter
what choices, the composition always equals the identity
map on $\det(L^\bullet)$. This clearly proves independence of those choices.

\medskip\noindent
The statement on canonical elements follows immediately from
Lemma \ref{lemma-canonical-element-well-defined} and our construction.
\end{proof}

\end{document}
